% =====================================================================
%  One-page project brief — Klimabonus field experiment
%  Purpose: circulate to colleagues for feedback. The single ask:
%  "What is the ONE question you would most want answered in a sample
%   and design like this?"
%
%  Compiles with pdflatex (TeX Live / Overleaf). Standard packages only.
% =====================================================================
\documentclass[10pt,a4paper]{article}

\usepackage[utf8]{inputenc}
\usepackage[T1]{fontenc}
\usepackage{lmodern}
\usepackage[english]{babel}
\usepackage{microtype}
\usepackage[a4paper,top=1.8cm,bottom=1.6cm,left=2.0cm,right=2.0cm]{geometry}
\usepackage{xcolor}
\usepackage{enumitem}
\usepackage{titlesec}
\usepackage{eurosym}
\usepackage[colorlinks=true,linkcolor=black,urlcolor=blue]{hyperref}

\definecolor{gublue}{HTML}{003C8F}
\definecolor{gugrey}{HTML}{555555}

% Compact section headings
\titleformat{\section}{\sffamily\bfseries\color{gublue}\large}{}{0pt}{}
\titlespacing*{\section}{0pt}{6pt}{2pt}

% Tight lists
\setlist[itemize]{leftmargin=1.3em,topsep=1pt,itemsep=1pt,parsep=0pt}

\setlength{\parindent}{0pt}
\setlength{\parskip}{3pt}
\pagestyle{empty}
\sloppy

\newcommand{\heading}[1]{{\sffamily\bfseries\color{gublue}#1}}

\begin{document}

% ---------- Title block ----------
{\sffamily\bfseries\LARGE\color{gublue}
Information Frictions and the Source of Information\\
in Firms' Take-up of Environmental Subsidies}

\vspace{2pt}
{\sffamily\color{gugrey}A randomised field experiment with the City of Frankfurt}

\vspace{4pt}
{\small\color{gugrey}%
[Author list — e.g. Guido Friebel, Jakob Famulok (Goethe University Frankfurt) and ZEW Mannheim] \quad\textbullet\quad Draft for feedback, \today}

\vspace{2pt}
{\color{gublue}\hrule height 1pt}
\vspace{4pt}

% ---------- Setting ----------
\section*{Setting}
The City of Frankfurt runs ``Klimabonus,'' a subsidy program co-funding
firms' investments in energy-reducing and climate-adaptation measures
(e.g.\ solar, green roofs, battery storage, fa\c{c}ade greening, rainwater
storage). Despite attractive terms, take-up among local firms is well
below administrative targets. We run the study in direct cooperation with
the city's climate office (\textit{Klimareferat}), which provides the
official program statistics and the administrative outcome data.

% ---------- Goal ----------
\section*{Goal of the study}
We ask \textbf{why firm take-up of a generous environmental subsidy
remains low}, and in particular whether the binding constraint is an
\textbf{information friction} and how it interacts with the \textbf{source
of information}. The design separates (i) the mere \emph{provision} of
credible program statistics from (ii) the \emph{relatability / source} of
that information, delivered through the experience of a comparable peer
firm. Beyond the average effect, we are interested in how prior beliefs
(success probability, time cost, expected award) move and whether belief
updating predicts real take-up.

% ---------- Treatment ----------
\section*{Treatment}
Contacted firms are randomised at the firm level into three arms that
\textbf{hold core program content constant} and vary only the surrounding
information:
\begin{itemize}
  \item \heading{Control:} a generic, neutral description of the program.
  \item \heading{T1 -- Statistics:} the same description plus official
        city statistics --- a \textbf{96\,\% approval rate}, \textbf{$\sim$25
        minutes} to complete an application, and an \textbf{average award
        of $\sim$\,\euro{}8{,}800} per measure.
  \item \heading{T2 -- Statistics + Peer:} the identical statistics plus a
        short testimonial (photo + quote) from a \textbf{size-matched peer
        firm} that successfully used the program.
\end{itemize}
The numbers are constant across T1 and T2; only the peer element is added
in T2, isolating the role of a relatable, same-type source.

% ---------- Sample & Flow ----------
\section*{Sample \& Flow}
\begin{itemize}
  \item \heading{Population:} $\sim$\,15{,}000 Frankfurt firms, identified
        from a firm panel and stratified by industry $\times$ size.
  \item \heading{Pure control:} a randomly selected group receives
        \textbf{no outreach at all} and is observed only through
        administrative records. It benchmarks the natural baseline take-up
        and identifies the effect of being contacted at all. All remaining
        (``contacted'') firms are randomised into the three arms above.
  \item \heading{Outreach:} addressed to the firm head; first by
        personalised business e-mail, then by postal letter to
        non-respondents. Each firm is directed to a treatment-specific
        landing page via a personalised link.
  \item \heading{Elicitation:} a short, voluntary survey on the landing
        page (see \textit{Outcomes}).
\end{itemize}

% ---------- Outcomes ----------
\section*{Outcomes}
\textbf{Primary:} actual engagement with the program --- application
submission, recorded in the city's administrative system over a 12-month
window and linked pseudonymously to the assigned arm.

\textbf{Secondary --- in survey.} Posterior beliefs, elicited as the key
mechanism: (i) the perceived \emph{approval probability}; (ii) the
\emph{expected award} (\euro); (iii) the \emph{time and effort} to file an
application and to claim the payout; and (iv) the expected \emph{processing
time} and overall perceived \emph{burden}. We further record stated
\emph{application intention} (0--10), the concrete \emph{measures} a firm
considers, and the \emph{barriers} it perceives, plus the stated importance
of climate action as a moderator.

\textbf{Secondary --- behavioural.} Click-through to the application portal
(tracked live on the landing page) and downstream technology adoption where
observable.

% ---------- The ask ----------
\section*{We would value your feedback}
If you could have us answer \textbf{one question} with a sample and design
like this, \textbf{what would it be?} We are especially keen on questions we
could still build in --- additional outcomes or belief items, a mechanism we
should isolate, a moderator or sub-sample worth powering for, or a threat to
identification we should pre-empt.

\end{document}
